\pdfbookmark{Общая характеристика работы}{characteristic}
\section*{Общая характеристика работы}

\newcommand{\actuality}{\pdfbookmark[1]{Актуальность}{actuality}\underline{\textbf{\actualityTXT}}}
\newcommand{\progress}{\pdfbookmark[1]{Разработанность темы}{progress}\underline{\textbf{\progressTXT}}}
\newcommand{\aim}{\pdfbookmark[1]{Цели}{aim}\underline{{\textbf\aimTXT}}}
\newcommand{\tasks}{\pdfbookmark[1]{Задачи}{tasks}\underline{\textbf{\tasksTXT}}}
\newcommand{\aimtasks}{\pdfbookmark[1]{Цели и задачи}{aimtasks}\aimtasksTXT}
\newcommand{\novelty}{\pdfbookmark[1]{Научная новизна}{novelty}\underline{\textbf{\noveltyTXT}}}
\newcommand{\influence}{\pdfbookmark[1]{Практическая значимость}{influence}\underline{\textbf{\influenceTXT}}}
\newcommand{\methods}{\pdfbookmark[1]{Методология и методы исследования}{methods}\underline{\textbf{\methodsTXT}}}
\newcommand{\defpositions}{\pdfbookmark[1]{Положения, выносимые на защиту}{defpositions}\underline{\textbf{\defpositionsTXT}}}
\newcommand{\reliability}{\pdfbookmark[1]{Достоверность}{reliability}\underline{\textbf{\reliabilityTXT}}}
\newcommand{\probation}{\pdfbookmark[1]{Апробация}{probation}\underline{\textbf{\probationTXT}}}
\newcommand{\contribution}{\pdfbookmark[1]{Личный вклад}{contribution}\underline{\textbf{\contributionTXT}}}
\newcommand{\publications}{\pdfbookmark[1]{Публикации}{publications}\underline{\textbf{\publicationsTXT}}}

\makeatletter
\@ifundefined{c@totalchapter@totc}{\newcounter{totalchapter@totc}}{}
\@ifundefined{c@totalappendix@totc}{\newcounter{totalappendix@totc}}{}
\@ifundefined{c@TotPages@totc}{\newcounter{TotPages@totc}}{}
\@ifundefined{c@totalcount@figure@totc}{\newcounter{totalcount@figure@totc}}{}
\@ifundefined{c@totalcount@table@totc}{\newcounter{totalcount@table@totc}}{}
\@ifundefined{c@citenum@totc}{\newcounter{citenum@totc}}{}
\setcounter{totalchapter@totc}{7}
\setcounter{totalappendix@totc}{2}
\setcounter{TotPages@totc}{142}
\setcounter{totalcount@figure@totc}{17}
\setcounter{totalcount@table@totc}{13}
\setcounter{citenum@totc}{70}
\makeatother


\textbf{Актуальность темы} определяется тем, что для современных edge/IoT-сред
одновременно характерны высокая динамичность, распределённый характер
управления и ограниченность локальных вычислительных ресурсов. В таких
условиях методы, основанные только на детерминированной оптимизации, требуют
слишком сильных предположений о полноте и стабильности информации о системе,
тогда как purely learning-based подходы, в том числе RL и MARL, хотя и обладают
адаптивностью, обычно не содержат встроенного механизма экономически
интерпретируемого согласования интересов участников.

Отдельную трудность представляет то, что в распределённой edge-среде решения
принимаются многократно и последовательно: текущее назначение задач влияет на
будущие очереди, загрузку узлов, сетевые задержки и последующие возможности
распределения ресурсов. Следовательно, задача должна рассматриваться не как
одноразовая оптимизация, а как последовательный процесс управления в
многоагентной частично наблюдаемой среде.

При этом в литературе сохраняется разрыв между двумя линиями исследований. С одной стороны, \textbf{аукционные и механизм-дизайн подходы} позволяют формально описывать стимулы, платежи, правдивость заявок и максимизацию общественного благосостояния, однако в классическом виде хуже приспособлены к динамической нестационарной среде. С другой стороны, \textbf{кооперативное многоагентное обучение с подкреплением} позволяет строить адаптивные политики в сложной среде, но само по себе не обеспечивает явного механизма стимулирования и интерпретируемого учёта внешних эффектов перераспределения ресурсов.

Тем самым актуальной является задача синтеза такого подхода, который сочетал бы экономически осмысленный механизм распределения, строгую многоагентную формализацию и адаптивную обучаемую политику, пригодную для децентрализованного исполнения в edge-среде.

\textbf{Целью} настоящей диссертационной работы является разработка моделей и гибридного метода адаптивного распределения вычислительных ресурсов в распределённых многоагентных edge-инфраструктурах в условиях частичной наблюдаемости и динамической нагрузки на основе интеграции аукционного механизма VCG и кооперативного многоагентного обучения с подкреплением.

Для достижения поставленной цели необходимо решить следующие \textbf{задачи}.
\begin{enumerate}
	\item Провести анализ существующих методов распределения ресурсов в edge- и IoT-системах, выявить ограничения детерминированных, аукционных и RL/MARL-подходов и обосновать требования к разрабатываемому методу;
	\item Разработать математическую модель распределения ресурсов в распределённой edge-системе с учётом гетерогенности ресурсов, ограничений по задержке, характеристик задач и экономических параметров взаимодействия агентов;
	\item Формализовать задачу адаптивного управления ресурсами в виде расширенной модели Dec-POMDP и обосновать применение кооперативного MARL-подхода в парадигме централизованного обучения и децентрализованного исполнения;
	\item Разработать гибридный метод и алгоритм MA--VCG--QMIX, интегрирующий аукционный механизм распределения ресурсов и обучаемую многоагентную политику принятия решений;
	\item Провести экспериментальную оценку предложенного метода в ряде сценариев и определить его преимущества, ограничения и области применимости.
	\item Определить условия практического применения, ограничения, архитектурное место и переносимость предложенного метода для прикладных edge-сценариев, включая роботизированные складские системы, промышленный AIoT и интеллектуальные транспортные подсистемы.
\end{enumerate}

\textbf{Научная новизна} диссертации состоит в следующем:
\begin{enumerate}
	\item Предложена расширенная модель распределения вычислительных ресурсов в edge-системе, объединяющая ресурсные ограничения, сетевые задержки, характеристики задач и экономические параметры взаимодействия агентов;
	\item Предложена формализация задачи адаптивного распределения ресурсов в виде расширенной модели Dec-POMDP, в которой экономические сигналы аукционного механизма включены в контур принятия решений агентов;
	\item Разработан гибридный метод MA--VCG--QMIX, сочетающий VCG-подобный аукционный слой и кооперативное многоагентное обучение с подкреплением в парадигме CTDE;
	\item Предложен способ включения аукционных сигналов в пространство состояний, наблюдений, допустимых действий и вознаграждения обучаемой многоагентной системы;
	\item Разработана программа экспериментальной валидации предложенного метода, обеспечивающая проверку его эффективности, адаптивности, устойчивости и fairness-свойств на воспроизводимом наборе сценариев.
	\item Экспериментально установлены условия применимости предложенного гибридного подхода в сценариях стабильной, пиковой и нарушенной нагрузки, а также выявлен компромисс между эффективностью распределения, справедливостью и адаптивностью.
\end{enumerate}

\textbf{Теоретическая значимость} работы состоит в развитии формализации задач распределения ресурсов в многоагентных edge-системах с частичной наблюдаемостью, а также в исследовании возможностей интеграции механизмов стимулирования и кооперативного MARL в едином контуре управления.

\textbf{Практическая значимость} работы состоит в возможности применения разработанного метода для построения систем распределения вычислительных задач в edge-инфраструктурах с динамической нагрузкой и ограниченной наблюдаемостью, в частности в роботизированных складских системах, промышленном AIoT, распределённой робототехнике и интеллектуальных транспортных подсистемах.

\textbf{Обоснованность и достоверность} полученных результатов обеспечивается:
\begin{enumerate}
	\item использованием признанных теоретических оснований: теории игр, аукционной теории, теории марковских процессов принятия решений, Dec-POMDP и методов MARL;
	\item согласованностью математической модели распределения ресурсов, формальной многоагентной постановки и гибридного алгоритмического решения;
	\item воспроизводимым протоколом экспериментальной валидации, включающим сценарное тестирование, сравнение с baseline-методами, многократные независимые запуски и статистическую обработку результатов;
	\item возможностью независимой проверки алгоритма на реализации в репозитории с использованием фиксированных конфигураций среды и параметров обучения.
\end{enumerate}

\textbf{Апробация работы.} Результаты, изложенные в данной диссертации, докладывались и обсуждались на конференциях:
\begin{enumerate}
	\item 11-я Международная конференция <<Распределенные вычисления и грид-технологии в науке и образовании>> (GRID'2025), Российская Федерация, г. Дубна, 07.07.2025–11.07.2025;
	\item Ежегодная всероссийская научная конференция по проблемам информатики в секции <<Методы машинного обучения и архитектуры интеллектуальных систем>>, Российская Федерация, г. Санкт-Петербург, 28.04.2026.
\end{enumerate}

\textbf{Публикации.} Основные результаты по теме диссертации изложены в 4 печатных изданиях, из которых ... — в периодических научных журналах, индексируемых Scopus [...], ... — в периодических научных журналах, индексируемых РИНЦ [...]. Получены 1 свидетельства о государственной регистрации программ для ЭВМ [...]. TODO 

\textbf{Основные положения, выносимые на защиту:}
\begin{enumerate}
	\item Математическая модель распределения вычислительных ресурсов в распределенных многоагентных системах периферийных вычислений, учитывающая гетерогенность узлов, ограничения по вычислительным ресурсам и задержкам, а также экономические характеристики взаимодействия агентов.
	\item Формализация задачи адаптивного управления распределением ресурсов в виде расширенной модели Dec-POMDP с частичной наблюдаемостью и стохастической динамикой среды.
	\item Гибридный метод MA-VCG-QMIX, объединяющий аукционное распределение ресурсов по механизму VCG и кооперативное многоагентное обучение с подкреплением в парадигме CTDE.
	\item Способ включения экономических сигналов аукционного механизма в контур обучения агентов, обеспечивающий совместный учет эффективности распределения, адаптивности и справедливости.
	\item Программа экспериментальной валидации и бенчмаркинга, предназначенная для воспроизводимой проверки эффективности, устойчивости и границ применимости предложенного метода.
	\item Результаты экспериментального исследования, показывающие, что в сценариях стабильной и пиковой нагрузки предложенный метод улучшает показатели распределения ресурсов по сравнению с базовыми подходами, а в сценариях высокой неоднородности и отказов выявляются границы его применимости.
\end{enumerate}

\textbf{Объем и структура работы.} Диссертация состоит из~введения,
\formbytotal{totalchapter}{глав}{ы}{}{},
заключения и
\formbytotal{totalappendix}{приложен}{ия}{ий}{}.
%% на случай ошибок оставляю исходный кусок на месте, закомментированным
%Полный объём диссертации составляет  \ref*{TotPages}~страницу
%с~\totalfigures{}~рисунками и~\totaltables{}~таблицами. Список литературы
%содержит \total{citenum}~наименований.
%
Полный объём диссертации составляет
\formbytotal{TotPages}{страниц}{у}{ы}{}, включая
\formbytotal{totalcount@figure}{рисун}{ок}{ка}{ков} и
\formbytotal{totalcount@table}{таблиц}{у}{ы}{}.
Список литературы содержит
\formbytotal{citenum}{наименован}{ие}{ия}{ий}.

\textbf{Во введении} сформулированы критерии, показана актуальность и новизна исследования, описана теоретическая и практическая значимость полученных результатов, обозначена цель и задачи исследования для ее достижения, сформулированы положения, выносимые на защиту.

\textbf{В первой главе} проведен анализ существующих подходов к управлению вычислительными ресурсами в распределенных системах периферийных вычислений. Рассмотрены детерминированные, аукционные и многоагентные методы, а также подходы на основе обучения с подкреплением. Выявлены их ограничения и обоснована необходимость разработки гибридного метода.

\textbf{Во второй главе} разработана математическая модель распределения ресурсов в распределенной edge-системе. Сформулированы модели задач и ресурсов, введены функции полезности и стоимости, а также формализована задача максимизации социального благосостояния и описан аукционный механизм распределения ресурсов.

\textbf{В третьей главе} предложена модель адаптивного управления ресурсами на основе многоагентного обучения с подкреплением. Задача формализована в виде Dec-POMDP, описаны состояния, наблюдения, действия, функция вознаграждения и обоснован выбор алгоритма QMIX.

\textbf{В четвертой главе} разработан гибридный метод управления ресурсами, объединяющий аукционный механизм VCG и алгоритм QMIX. Предложена расширенная модель Dec-POMDP, интегрирующая экономические сигналы, и сформулирован алгоритм MA--VCG--QMIX.

\textbf{В пятой глава} приведено сформулирована программа экспериментальной валидации и
бенчмаркинга предложенного метода. Программа включает основные элементы формальной модели: устройства, edge-узлы, сетевую топологию, поток задач, аукционное распределение, VCG-платежи, reward-сигнал и обучение координированной политики.

\textbf{В шестой главе} проведена экспериментальная оценка эффективности разработанного метода. Рассмотрены различные сценарии работы системы, выполнено сравнение с базовыми подходами и проанализированы полученные результаты.

\textbf{В седьмой главе } рассмотрены практические аспекты применения предложенного метода, его место в прикладной архитектуре, условия внедрения, ограничения и переносимость на роботизированные складские системы, промышленный AIoT, интеллектуальные транспортные подсистемы и другие edge-сценарии.

\textbf{В заключении} подведены итоги исследования, обсуждены ограничения, определены направления дальнейшей работы.


\pdfbookmark{Содержание работы}{description}
\section*{Содержание работы}

Диссертация посвящена разработке гибридного метода адаптивного распределения
вычислительных задач автономных мобильных роботов между периферийными
вычислительными узлами. Логика исследования выстроена от анализа предметной
области и формальной постановки задачи к построению аукционного и обучаемого
контуров, синтезу полного метода, созданию программного прототипа и его
экспериментальной проверке в воспроизводимой складской AMR-среде.

\underline{\textbf{Во введении}} обоснована актуальность темы, сформулированы
цель и задачи исследования, определены объект и предмет исследования,
показаны научная новизна, теоретическая и практическая значимость работы,
сформулированы положения, выносимые на защиту, приведены сведения об
апробации и публикациях автора.

\underline{\textbf{В первой главе}} проведён анализ предметной области
автономных мобильных роботов с поддержкой периферийных вычислений и
существующих подходов к распределению вычислительных ресурсов. Показано, что
детерминированные оптимизационные методы, аукционные механизмы и методы
обучения с подкреплением по отдельности закрывают только часть требований,
возникающих в локальных AMR-системах. На этой основе обоснована необходимость
гибридной архитектуры, сочетающей экономически интерпретируемое распределение
дефицитных ресурсов с адаптивным многоагентным управлением периферийными
узлами.

\underline{\textbf{Во второй главе}} разработана доменная и математическая
модель AMR-системы с поддержкой периферийных вычислений. В модели заданы
множества роботов, периферийных узлов и вычислительных задач, граф
инфраструктуры, классы задач, параметры сетевой задержки и очередей,
функции полезности роботов и функции стоимости использования ресурсов.
Сформулирован критерий общественного благосостояния, а также показано, как
операционные режимы узлов задают ограничения для последующего аукционного
раунда распределения задач.

\underline{\textbf{В третьей главе}} построен аукционный слой распределения
задач AMR между периферийными узлами на основе VCG-подобного механизма.
Определены структура заявок, задача выбора победителей, правила вычисления
платежей и интерпретация внешних эффектов использования дефицитных ресурсов.
Отдельно зафиксированы границы применимости теоретических свойств: локальная
эффективность, стратегическая совместимость и индивидуальная рациональность
рассматриваются только для отдельного статического раунда распределения.
Основные результаты главы отражены в статье~\cite{Tereschenko2026MAVCG}.

\underline{\textbf{В четвёртой главе}} задача адаптивного управления
периферийными узлами формализована как децентрализованный частично
наблюдаемый марковский процесс принятия решений. Обучаемыми агентами
выступают периферийные узлы, а их действиями являются операционные режимы
допуска, резервирования и обслуживания задач. Обосновано применение
парадигмы централизованного обучения и децентрализованного исполнения, а в
качестве базового алгоритма кооперативного обучения выбран QMIX, согласующийся
с дискретным пространством режимов, частичной наблюдаемостью и командным
вознаграждением.

\underline{\textbf{В пятой главе}} предложен гибридный метод AMPLE-AMR,
объединяющий обучаемую многоагентную политику выбора операционных режимов
периферийных узлов и VCG-подобный аукционный слой распределения задач.
Показано, что обучаемый контур не формирует ставки и не изменяет оценки
ценности задач, а отвечает только за адаптивную настройку условий
распределения. В той же главе предложен масштабируемый вариант C-AMPLE-AMR,
основанный на кластеризации графа периферийной инфраструктуры, локальных
аукционах и критерии перехода от глобального режима распределения к
кластерному.

\underline{\textbf{В шестой главе}} описан программный прототип,
реализующий полный контур метода: генерацию задач AMR, моделирование
периферийной инфраструктуры и сети, эвристическое и VCG-подобное
распределение, многоагентное обучение режимов узлов, а также экспорт
CSV-результатов, графиков и \LaTeX-таблиц. Подчёркнуто, что в программной
реализации сохранено принципиальное разделение между экономическим
механизмом назначения задач и обучаемым контуром управления режимами узлов.
Разработанный прототип зарегистрирован как программа для
ЭВМ~\cite{patent_ma_vcg}.

\underline{\textbf{В седьмой главе}} выполнена экспериментальная валидация
методов AMPLE-AMR и C-AMPLE-AMR в параметризованной складской AMR-среде.
Проверка проводилась в сценариях стабильной нагрузки, пикового потока задач,
разнородной инфраструктуры, деградации связи, отказов периферийных узлов и
роста масштаба системы. Показано, что в сценарии стабильной нагрузки
AMPLE-AMR увеличивает общественное благосостояние на $12.6\%$ по сравнению с
Fixed-Heuristic, повышает долю завершённых задач с $0.693$ до $0.740$ и
снижает 95-й процентиль задержки примерно на $29.8\%$. В сценарии пиковых
всплесков нагрузки выигрыш по общественному благосостоянию составляет
$12.8\%$, а при деградации связи и отказах узлов гибридные методы устойчиво
превосходят фиксированные baselines по общественному благосостоянию и
задержке. Для конфигурации Warehouse-XL кластерная версия C-AMPLE-AMR
снижает накладные расходы принятия решения на $44.6\%$ при потере
общественного благосостояния около $1.57\%$, что подтверждает
целесообразность перехода к кластерному режиму только на достаточно крупных
масштабах.

\pdfbookmark{Заключение}{conclusion}
\section*{Заключение}

В заключении приведены основные результаты исследования:
%% Согласно ГОСТ Р 7.0.11-2011:
%% 5.3.3 В заключении диссертации излагают итоги выполненного исследования, рекомендации, перспективы дальнейшей разработки темы.
%% 9.2.3 В заключении автореферата диссертации излагают итоги данного исследования, рекомендации и перспективы дальнейшей разработки темы.

Во-первых, выполнен анализ предметной области автономных мобильных роботов с поддержкой периферийных вычислений. Показано, что существующие подходы к распределению ресурсов в периферийных системах развиваются по нескольким направлениям: детерминированная оптимизация, аукционные механизмы, обучение с подкреплением, многоагентное обучение и гибридные схемы. Детерминированные методы позволяют строго задавать ограничения и целевые функции, но требуют достаточно полной и устойчивой информации о состоянии системы. Аукционные механизмы позволяют учитывать дефицит ресурсов и внешние эффекты распределения, но в классическом виде ориентированы на отдельный раунд. Методы обучения с подкреплением позволяют учитывать долгосрочные последствия решений, но сами по себе не задают экономически интерпретируемого механизма распределения дефицитных ресурсов.

Во-вторых, разработана доменная и математическая модель AMR-системы с поддержкой периферийных вычислений. В модели определены множества автономных мобильных роботов, периферийных вычислительных узлов, вычислительных задач, граф сетевой связности, параметры задержек, очередей, отказов и ресурсных ограничений. Отдельно учтены классы робототехнических задач, различающиеся по критичности, срокам обработки, объёму данных и требованиям к ресурсам. Предложенная модель позволяет описывать не только вычислительную сложность задачи, но и её прикладную значимость для робота и текущей миссии.

В-третьих, построен VCG-подобный аукционный слой распределения вычислительных задач. В рамках одного раунда распределения задача назначения формализована как задача максимизации общественного благосостояния с учётом полезности выполнения задач, стоимости использования ресурсов, сетевых задержек, очередей, отказов и ограничений по срокам обработки. Платежи аукционного слоя интерпретируются как оценки внешнего эффекта, создаваемого участниками при использовании дефицитных ресурсов. При этом в работе явно зафиксированы границы применимости VCG-подобных свойств: эффективность, индивидуальная рациональность и стратегическая совместимость рассматриваются только для отдельного статического раунда при стандартных предпосылках и не переносятся автоматически на всю повторяющуюся систему с обучением.

В-четвёртых, задача адаптивного управления периферийными узлами формализована как децентрализованный частично наблюдаемый марковский процесс принятия решений. В предложенной постановке обучаемыми агентами являются периферийные вычислительные узлы, а не роботы, формирующие заявки. Действиями агентов являются операционные режимы узлов: режимы допуска, резервирования, приоритизации и защиты от перегрузки. Такое разделение ролей принципиально важно: обучаемая политика не управляет ставками и не изменяет оценки ценности задач, а значит, не подменяет аукционный механизм стратегическим обучением заявок.

В-пятых, обосновано применение кооперативного многоагентного обучения с подкреплением в парадигме централизованного обучения и децентрализованного исполнения. В качестве базового алгоритма выбран QMIX, поскольку он согласуется с дискретным пространством операционных режимов, командным вознаграждением, частичной наблюдаемостью и необходимостью локального исполнения политики периферийными узлами. Централизованное обучение может выполняться в серверном контуре, облаке, цифровом двойнике или имитационной среде, тогда как в режиме эксплуатации каждый узел выбирает режим на основе локального наблюдения.

В-шестых, разработан гибридный метод AMPLE-AMR. Метод объединяет обучаемую многоагентную политику выбора операционных режимов периферийных узлов и VCG-подобный аукционный слой распределения задач. В AMPLE-AMR обучаемая политика определяет условия распределения, то есть экспонируемую ёмкость и режимы обслуживания узлов, а аукционный слой выполняет назначение задач в рамках выбранных ресурсных ограничений. Такое построение позволяет совместить адаптацию к динамике среды и экономически интерпретируемое распределение дефицитных ресурсов.

В-седьмых, предложена масштабируемая версия C-AMPLE-AMR, основанная на кластеризации графа периферийной инфраструктуры и проведении локальных распределений внутри кластеров. Кластерная версия предназначена для случаев, когда глобальный аукционный раунд становится слишком дорогим по времени вычислений или объёму коммуникации. В работе введены критерии оценки целесообразности перехода к кластерному режиму, учитывающие накладные расходы глобального распределения, длительность управляющего слота, межкластерный срез графа и возможную потерю общественного благосостояния.

В-восьмых, реализован программный прототип AMPLE-AMR и сформирована воспроизводимая экспериментальная среда. Прототип включает генератор складских сценариев, модели роботов и периферийных узлов, эвристическое распределение, VCG-подобное распределение, кластерное VCG-подобное распределение, контур обучения режимов узлов и автоматический экспорт результатов в виде таблиц и графиков.

В-девятых, проведена экспериментальная проверка предложенного метода в сценариях стабильной нагрузки, пикового потока задач, разнородной периферийной инфраструктуры, деградации связи, отказов узлов, масштабирования и анализа чувствительности параметров режимов. Эксперименты показали, что наиболее устойчивый практический вклад вносит обучаемый выбор операционных режимов периферийных узлов. AMPLE-AMR превосходит Fixed-Heuristic по общественному благосостоянию во всех основных сценариях: на $124.0$ условных единицы при стабильной нагрузке, на $210.5$ при пиковой нагрузке, на $95.6$ в разнородной инфраструктуре, на $94.9$ при деградации связи и на $110.7$ при отказах узлов. При этом самостоятельное преимущество VCG-подобного распределения над жадной эвристикой при фиксированных режимах в текущей параметризации не подтверждено, а диагностическая роль внутренних цен подтверждена лишь частично.


\pdfbookmark{Литература}{bibliography}
\nocite{Tereschenko2026MAVCG}
\nocite{patent_ma_vcg}
\ifdefmacro{\microtypesetup}{\microtypesetup{protrusion=false}}{}
\urlstyle{rm}
\ifnumequal{\value{bibliosel}}{0}{%
    \renewcommand{\bibname}{\large Литература и публикации автора}
    \renewcommand{\refname}{\large Литература и публикации автора}
    \insertbibliofull
}{%
    \insertbiblioauthor
    \begin{refcontext}[labelprefix={}]
        \insertbiblioexternal
    \end{refcontext}
}
\ifdefmacro{\microtypesetup}{\microtypesetup{protrusion=true}}{}
\urlstyle{tt}
